\Needspace{8\baselineskip}
\Article{273}
Hier folgen einige Beispiele, die die vorigen Regeln näher erläutern.

\emph{Beispiel} 1. Es sei \(f=\left(\begin{smallmatrix}19,&21,&50\\15,&28,&1\end{smallmatrix}\right)\); dann gilt \(F=\left(\begin{smallmatrix}-825,&-166,&-398\\257,&573,&-370\end{smallmatrix}\right)\), \(D=-1\). Da \((19,1,21)\) eine reduzierte binäre Form ist, zu der keine andere mit einem ersten Koeffizienten kleiner als \(19\) äquivalent ist, lässt sich die erste Reduktion hier nicht anwenden. Die Theorie der Äquivalenz binärer Formen zeigt aber, dass die binäre Form
\SourceLeaf{316}
\((A'',B,A')=(-398,257,-166)\) in die einfachere äquivalente Form \((-2,1,-10)\) verwandelt werden kann, in die sie durch die Substitution \(2,7,3,11\) übergeht. Setzt man also \(\beta'=2\), \(\gamma'=-7\), \(\beta''=-3\), \(\gamma''=11\), so muss auf die Form \(f\) die Substitution \(\left\{\begin{smallmatrix}1,&0,&0\\0,&2,&-7\\0,&-3,&11\end{smallmatrix}\right\}\) angewendet werden, durch die sie in \(\left(\begin{smallmatrix}19,&354,&4769\\-1299,&301,&-82\end{smallmatrix}\right)\ldots f'\) übergeht. Der dritte Koeffizient der dieser Form adjungierten Form ist \(-2\); in dieser Hinsicht ist also \(f'\) als einfacher als \(f\) anzusehen.

Auf die Form \(f'\) lässt sich die erste Reduktion anwenden. Da nämlich die binäre Form \((19,-82,354)\) in \((1,0,2)\) durch die Substitution \(13,4,3,1\) verwandelt wird: muss auf die Form \(f'\) die Substitution \(\left\{\begin{smallmatrix}13,&4,&0\\3,&1,&0\\0,&0,&1\end{smallmatrix}\right\}\) angewendet werden, durch die sie in \(\left(\begin{smallmatrix}1,&2,&4769\\-95,&16,&0\end{smallmatrix}\right)\ldots f''\) übergeht.

Auf die Form \(f''\), der die Form \(\left(\begin{smallmatrix}-513,&-4513,&-2\\-95,&32,&1520\end{smallmatrix}\right)\) adjungiert ist, lässt sich wieder die zweite Reduktion anwenden. Denn \((-2,-95,-4513)\) geht durch die Substitution \(47,1,-1,0\) in \((-1,1,-2)\) über: deshalb muss auf \(f''\) die Substitution \(\left\{\begin{smallmatrix}1,&0,&0\\0,&47,&-1\\0,&1,&0\end{smallmatrix}\right\}\) angewendet werden, durch die sie in \(\left(\begin{smallmatrix}1,&257,&2\\1,&0,&16\end{smallmatrix}\right)\ldots f'''\) übergeht. Deren erster Koeffizient kann durch die erste Reduktion nicht weiter verkleinert werden, ebenso wenig der dritte Koeffizient der ihr adjungierten Form durch die zweite.

\emph{Beispiel} 2. Es sei die Form \(\left(\begin{smallmatrix}10,&26,&2\\7,&0,&4\end{smallmatrix}\right)\ldots f\) gegeben, der die Form \(\left(\begin{smallmatrix}-3,&-20,&-244\\70,&-28,&8\end{smallmatrix}\right)\) adjungiert ist und deren Determinante \(=2\). Wendet man abwechselnd die zweite und die erste Reduktion an, so erhält man nacheinander
\[
\begin{array}{c|c|c}
\text{Substitutionen}&\text{überführen}&\text{in}\\[3pt]
\left\{\begin{smallmatrix}1,&0,&0\\0,&-1,&0\\0,&4,&-1\end{smallmatrix}\right\}
&f&\left(\begin{smallmatrix}10,&2,&2\\-1,&0,&-4\end{smallmatrix}\right)=f'\\[6pt]
\left\{\begin{smallmatrix}0,&-1,&0\\1,&-2,&0\\0,&0,&1\end{smallmatrix}\right\}
&f'&\left(\begin{smallmatrix}2,&2,&2\\2,&-1,&0\end{smallmatrix}\right)=f''\\[6pt]
\left\{\begin{smallmatrix}1,&0,&0\\0,&-1,&0\\0,&2,&-1\end{smallmatrix}\right\}
&f''&\left(\begin{smallmatrix}2,&2,&2\\-2,&1,&-2\end{smallmatrix}\right)=f'''\\[6pt]
\left\{\begin{smallmatrix}1,&0,&0\\1,&1,&0\\0,&0,&1\end{smallmatrix}\right\}
&f'''&\left(\begin{smallmatrix}0,&2,&2\\-2,&-1,&0\end{smallmatrix}\right)=f''''
\end{array}
\]
Die Form \(f''''\) lässt sich durch die erste oder zweite Reduktion nicht weiter herabsetzen.

\Needspace{8\baselineskip}
\Article{274}
Hat man eine ternäre Form, deren erster Koeffizient sowie der dritte Koeffizient ihrer adjungierten Form durch die vorigen Methoden so weit wie möglich herabgesetzt worden sind: so liefert die folgende Methode eine weitere Reduktion.

\SourceLeaf{317}
Verwendet man dieselben Bezeichnungen wie in Artikel 272 und setzt \(\alpha=1\), \(\alpha'=0\), \(\beta'=1\), \(\alpha''=0\), \(\beta''=0\), \(\gamma''=1\), \emph{das heißt}, verwendet man die Substitution
\[
\begin{array}{ccc}
1,&\beta,&\gamma\\
0,&1,&\gamma'\\
0,&0,&1
\end{array}
\]
so gilt
\[
\begin{aligned}
m&=a,\quad m'=a'+2b''\beta+a\beta\beta,\\
m''&=a''+2b\gamma'+2b'\gamma+a\gamma\gamma+2b''\gamma\gamma'+a'\gamma'\gamma'\\
n&=b+a'\gamma'+b'\beta+b''(\gamma+\beta\gamma')+a\beta\gamma,\\
n'&=b'+a\gamma+b''\gamma',\quad n''=b''+a\beta
\end{aligned}
\]
und außerdem
\[
M''=A'',\quad N=B-A''\gamma',\quad N'=B'-N\beta-A''\gamma
\]
Durch eine solche Substitution werden also die Koeffizienten \(a\), \(A''\), die durch die vorigen Reduktionen verkleinert worden sind, nicht verändert. Die Aufgabe besteht daher darin, durch geeignete Wahl von \(\beta\), \(\gamma\), \(\gamma'\) die übrigen Koeffizienten zu verkleinern. Dazu bemerken wir zunächst: Ist \(A''=0\), so darf man annehmen, dass auch \(a=0\) gilt. Denn wäre \(a\) nicht \(=0\), so ließe sich die erste Reduktion noch einmal anwenden, da jede binäre Form mit der Determinante \(0\) zu einer Form der Art \((0,0,h)\) äquivalent ist, also zu einer Form, deren erster Koeffizient \(=0\) ist (siehe Artikel 215). Auf genau dieselbe Weise darf man annehmen, dass auch \(A''=0\) gilt, wenn \(a=0\), sodass entweder keine der Zahlen \(a\), \(A''\) gleich \(0\) ist oder beide.

Im ersten Fall ist klar, dass \(\beta\), \(\gamma\), \(\gamma'\) so gewählt werden können, dass ohne Rücksicht auf das Vorzeichen \(n''\), \(N\), \(N'\) jeweils nicht größer als \(\frac{1}{2}a\), \(\frac{1}{2}A''\), \(\frac{1}{2}A''\) sind. So geht im ersten Beispiel des vorigen Artikels die letzte Form \(\left(\begin{smallmatrix}1,&257,&2\\1,&0,&16\end{smallmatrix}\right)\), der die Form \(\left(\begin{smallmatrix}-513,&-2,&-1\\1,&-16,&32\end{smallmatrix}\right)\) adjungiert ist, durch die Substitution \(\left\{\begin{smallmatrix}1,&-16,&16\\0,&1,&-1\\0,&0,&1\end{smallmatrix}\right\}\) in \(\left(\begin{smallmatrix}1,&1,&1\\0,&0,&0\end{smallmatrix}\right)\ldots f''''\) über, der die Form \(\left(\begin{smallmatrix}-1,&-1,&-1\\0,&0,&0\end{smallmatrix}\right)\) adjungiert ist.

Im zweiten Fall, in dem \(a=A''=0\) und damit auch \(b''=0\) gilt, erhält man
\[
\begin{aligned}
m&=0,\quad m'=a',\quad m''=a''+2b\gamma'+2b'\gamma+a'\gamma'\gamma'\\
n&=b+a'\gamma'+b'\beta,\quad n'=b',\quad n''=0
\end{aligned}
\]
Also gilt
\[
D=a'b'b'=m'n'n'
\]
und man sieht leicht, dass \(\beta\) und \(\gamma'\) so bestimmt werden können, dass \(n\) gleich dem absolut kleinsten Rest von \(b\) nach dem Modul wird, der der größte gemeinsame Teiler
\SourceLeaf{318}
von \(a'\), \(b'\) ist; das heißt, dass \(n\) ohne Rücksicht auf das Vorzeichen nicht größer als die Hälfte dieses Teilers wird, also \(n=0\), sobald \(a'\), \(b'\) teilerfremd sind. Sind \(\beta\), \(\gamma'\) auf diese Weise bestimmt, so lässt sich der Wert von \(\gamma\) so wählen, dass \(m''\) ohne Rücksicht auf das Vorzeichen nicht größer als \(b'\) ist. Dies wäre zwar unmöglich, wenn \(b'=0\); dann wäre aber \(D=0\), welchen Fall wir ausgeschlossen haben. So erhält man für die letzte Form in \emph{Beispiel} 2 des vorigen Artikels \(n=-2-\beta+2\gamma'\), woraus man durch \(\beta=-2\), \(\gamma'=0\) den Wert \(n=0\) erhält. Ferner gilt \(m''=2-2\gamma\), und mit \(\gamma=1\) erhält man \(m''=0\). Wir haben also die Substitution \(\left\{\begin{smallmatrix}1,&-2,&1\\0,&1,&0\\0,&0,&1\end{smallmatrix}\right\}\), durch die jene Form in \(\left(\begin{smallmatrix}0,&2,&0\\0,&-1,&0\end{smallmatrix}\right)\ldots f'''''\) übergeht.

\Needspace{8\baselineskip}
\Article{275}
Hat man eine Folge äquivalenter ternärer Formen \(f\), \(f'\), \(f''\), \(f'''\) usw.\ und Transformationen jeder dieser Formen in die folgende: so lässt sich aus den Transformationen der Form \(f\) in \(f'\) und der Form \(f'\) in \(f''\) nach Artikel 270 eine Transformation der Form \(f\) in \(f''\) herleiten. Aus dieser und der Transformation der Form \(f''\) in \(f'''\) folgt eine Transformation der Form \(f\) in \(f'''\) usw.\ Offenbar lässt sich so eine Transformation der Form \(f\) in jede andere Form der Folge finden. Da sich außerdem aus einer Transformation der Form \(f\) in eine beliebige andere äquivalente Form \(g\) eine Transformation der Form \(g\) in \(f\) herleiten lässt (\(S''\) aus \(S\), Artikel 268, 269), kann man auf diese Weise eine Transformation jeder Form der Folge \(f'\), \(f''\) usw.\ in die erste, \(f\), bestimmen. --- So findet man für die Formen des ersten Beispiels im vorigen Artikel die Substitutionen
\[
\begin{array}{rrr|rrr|rrr}
13,&4,&0&13,&188,&-4&13,&-20,&16\\
6,&2,&-7&6,&87,&-2&6,&-9,&7\\
-9,&-3,&11&-9,&-130,&3&-9,&14,&-11
\end{array}
\]
durch die \(f\) jeweils in \(f''\), \(f'''\), \(f''''\) übergeht, und aus der letzten Substitution diese \(\left\{\begin{smallmatrix}1,&4,&4\\3,&1,&5\\3,&-2,&3\end{smallmatrix}\right\}\), durch die \(f''''\) in \(f\) übergeht. Ebenso erhält man für \emph{Beispiel} 2 des vorigen Artikels die Substitutionen
\[
\begin{array}{rrr|rrr}
1,&-1,&1&2,&-3,&-1\\
-3,&4,&-3&3,&1,&0\\
10,&-14,&11&2,&4,&1
\end{array}
\]
durch die jeweils die Form \(\left(\begin{smallmatrix}10,&26,&2\\7,&0,&4\end{smallmatrix}\right)\) in \(\left(\begin{smallmatrix}0,&2,&0\\0,&-1,&0\end{smallmatrix}\right)\) und diese in jene übergeht.

\Needspace{8\baselineskip}
\Article{276}
\textsc{Satz.} \emph{Die Anzahl der Klassen, in die alle ternären Formen mit gegebener Determinante zerfallen, ist stets endlich.}

\SourceLeaf{319}
\emph{Beweis.} I. Die Anzahl aller Formen \(\left(\begin{smallmatrix}a,&a',&a''\\b,&b',&b''\end{smallmatrix}\right)\) mit gegebener Determinante \(D\), bei denen \(a=0\), \(b''=0\) gelten, \(b\) nicht größer als die Hälfte des größten gemeinsamen Teilers der Zahlen \(a'\), \(b'\) ist und \(a''\) nicht größer als \(b'\) ist, ist offenbar endlich. Da nämlich \(a'b'b'=D\) gelten muss, kann \(b'\) nur die Werte \(+1\), \(-1\) und die mit positivem und negativem Vorzeichen versehenen Wurzeln der Quadrate annehmen, die \(D\) teilen (falls es außer \(1\) noch andere gibt). Die Anzahl dieser Werte ist endlich. Für jeden Wert von \(b'\) ist aber der Wert von \(a'\) bestimmt, und die Werte von \(b\), \(a''\) sind offenbar auf eine endliche Anzahl beschränkt.

II. Ebenso ist die Anzahl aller Formen \(\left(\begin{smallmatrix}a,&a',&a''\\b,&b',&b''\end{smallmatrix}\right)\) mit der Determinante \(D\) endlich, bei denen \(a\) nicht \(=0\) und nicht größer als \(\frac{4}{3}\sqrt[3]{\pm D}\) ist; \(b''b''-aa'=A''\) nicht \(=0\) und nicht größer als \(\frac{4}{3}\sqrt[3]{D^2}\) ist; \(b''\) nicht größer als \(\frac{1}{2}a\) ist; und \(ab-b'b''=B\) und \(a'b'-bb''=B'\) nicht größer als \(\frac{1}{2}A''\) sind. Denn die Anzahl aller Kombinationen der Werte von \(a\), \(b''\), \(A''\), \(B\), \(B'\) ist endlich. Sind diese einzeln bestimmt, so sind auch die übrigen Koeffizienten der Form, \(a'\), \(b\), \(b'\), \(a''\), sowie die Koeffizienten der adjungierten Form
\[
bb-a'a''=A,\quad b'b'-aa''=A',\quad a''b''-bb'=B''
\]
durch folgende Gleichungen bestimmt:
\[
\begin{aligned}
a'&=\frac{b''b''-A''}{a},\quad A'=\frac{BB-aD}{A''},\\
A&=\frac{B'B'-a'D}{A''},\quad B''=\frac{BB'+b''D}{A''}\\
b&=\frac{AB-B'B''}{D}=-\frac{Ba'+B'b''}{A''},\\
b'&=\frac{A'B'-BB''}{D}=-\frac{Bb''+B'a}{A''}\\
a''&=\frac{b'b'-A'}{a}=\frac{bb-A}{a'}=\frac{bb'+B''}{b''}
\end{aligned}
\]
Nun erhält man alle jene Formen, indem man aus sämtlichen Kombinationen der Werte von \(a\), \(b''\), \(A''\), \(B\), \(B'\) diejenigen wählt, für die auch \(a'\), \(a''\), \(b\), \(b'\) ganzzahlige Werte annehmen; ihre Anzahl ist daher offenbar endlich.

III. Alle Formen in I und II bilden somit eine endliche Anzahl von Klassen. Diese kann sogar kleiner sein als die Anzahl der Formen selbst, wenn einige von ihnen einander äquivalent sind. Da nun nach den vorigen Untersuchungen jede ternäre Form mit der Determinante \(D\) einer jener Formen äquivalent sein muss, \emph{das heißt}, zu einer der von ihnen gebildeten Klassen gehören muss: umfassen diese Klassen alle Formen mit der Determinante \(D\), \emph{das heißt}, alle ternären Formen mit der Determinante \(D\) zerfallen in eine endliche Anzahl von Klassen. \emph{W.z.b.w.}
